\section{The circle's loop space: the 2013 encode--decode construction}
\label{sec:p07-hott-circle-encode-decode}

This section develops the main circle encode--decode construction of
\cite[\S8.1]{HoTT2013PublicReleaseFoundations},
at the public-release commit
\texttt{e0aa8e9d16858fc0ce01aacd43b24da0bd765ae3} of 19 June 2013.
It uses the circle eliminator of Definition~\ref{def:p07-hit-circle}, not a
judgmental computation rule for the loop. The path, transport, function-path,
univalence and equivalence interfaces are those already established in the
preceding HoTT foundations. Write \(p\cdot q\) for first traversing \(p\) and
then \(q\), and \(T_P(p,u)\) for transport. The notation \(=\) denotes an
identity type; \(\equiv_{\mathrm j}\) denotes judgmental equality.

\begin{hbconvention}{def:p07-circle-standing}{Universes and ordinary integers}
Work in the intensional dependent theory of
Definition~\ref{def:p07-hott-core}, with ordinary natural numbers
\(\mathbb N\), constructors \(0,\mathsf S\), their dependent eliminator and
judgmental constructor computations. Use finite sums and their dependent
eliminators. Funext and univalence are explicit assumptions, as in
Definitions~\ref{def:p07-hott-funext} and \ref{def:p07-hott-univalence}.
The circle has \(\mathsf{base}:S^1\) and
\(\mathsf{loop}:\mathsf{base}=\mathsf{base}\). A universe containing the
integers is used as a target of circle recursion in the next larger universe.
No UIP for \(S^1\), choice, resizing, or judgmental univalence is assumed.
\end{hbconvention}

\begin{hbdefinition}{def:p07-circle-integer-normalization}{The source's integer presentation}
Let \(C=1+\mathbb N+\mathbb N\), with tags
\(\mathsf z,\mathsf{pos}(k),\mathsf{neg}(k)\), representing \(0,k+1,-(k+1)\).
Define \(g:\mathbb N^2\to C\) by simultaneous cancellation:
\[
g(0,0)=\mathsf z,\quad g(\mathsf Sa,0)=\mathsf{pos}(a),\quad
g(0,\mathsf Sb)=\mathsf{neg}(b),\quad
g(\mathsf Sa,\mathsf Sb)=g(a,b).
\]
Let \(j(\mathsf z)=(0,0)\), \(j(\mathsf{pos}(a))=(\mathsf Sa,0)\),
\(j(\mathsf{neg}(b))=(0,\mathsf Sb)\), and \(r=jg\). We use
\[
\mathbb Z=\sum_{x:\mathbb N^2}(r(x)=x),\qquad
E(c)=(j(c),\sigma_c),\qquad G(x,h)=g(x),\qquad
q(x)=E(g(x)).
\]
Here \(\sigma_c:r(j(c))=j(c)\) is defined by tag induction, with
reflexivity in each constructor case. It is not asserted to be reflexivity
judgmentally for a free variable \(c\).
This is the fixed-point presentation of the idempotent canonicalizer in
\cite[\S6.10]{HoTT2013PublicReleaseFoundations}; the source's
relation on pairs is \((a,b)\sim(c,d)\) iff \(a+d=b+c\).
Put \(0=E(\mathsf z)\). Signed induction below means induction through
the equivalence \(E:C\simeq\mathbb Z\), not an extra axiom about the circle.
\end{hbdefinition}

\begin{hblemma}{thm:p07-circle-integer-presentation}{proved=annex:p07-circle-integer-presentation}{Canonical integers are the source quotient and are a set}
The type \(\mathbb N\) is a set, \(C\) and \(\mathbb Z\) are sets,
\(E\) and \(G\) are inverse equivalences, and
\[
g(a,b)=g(c,d)\quad\Longleftrightarrow\quad a+d=b+c.
\]
The map \(q\) has the universal property of the quotient by \(\sim\) into
sets. Thus the signed presentation does not replace the source integers by
an unrelated assumed type.
\end{hblemma}

\begin{hbdefinition}{def:p07-circle-integer-steps}{Successor, predecessor, and addition}
Define the following maps on \(C\):
\[
\begin{array}{c|cccc}
c&\mathsf z&\mathsf{pos}(k)&\mathsf{neg}(0)&\mathsf{neg}(\mathsf Sk)\\ \hline
s(c)&\mathsf{pos}(0)&\mathsf{pos}(\mathsf Sk)&\mathsf z&\mathsf{neg}(k)
\end{array}
\]
\[
\begin{array}{c|cccc}
c&\mathsf z&\mathsf{neg}(k)&\mathsf{pos}(0)&\mathsf{pos}(\mathsf Sk)\\ \hline
t(c)&\mathsf{neg}(0)&\mathsf{neg}(\mathsf Sk)&\mathsf z&\mathsf{pos}(k).
\end{array}
\]
Set \(\mathsf{succ}=EsG\), \(\mathsf{pred}=EtG\) on \(\mathbb Z\).
Integer addition is the descended pair operation
\[
q(a,b)+q(c,d)=q(a+c,b+d).
\]
All comparisons on arbitrary fixed-point pairs, including the preceding
formula, are propositional; they are not additional judgmental rules.
\end{hbdefinition}

\begin{hblemma}{thm:p07-circle-integer-steps}{proved=annex:p07-circle-integer-steps}{The successor equivalence and additive recurrences}
There is an equivalence \(e_{\mathsf{succ}}:\mathbb Z\simeq\mathbb Z\)
with forward map \(\mathsf{succ}\) and inverse \(\mathsf{pred}\).
Addition is well defined and satisfies
\[
m+0=m,\quad m+\mathsf{succ}(n)=\mathsf{succ}(m+n),\quad
m+\mathsf{pred}(n)=\mathsf{pred}(m+n).
\]
It also satisfies \(0+n=n\).
\end{hblemma}

\begin{hbdefinition}{def:p07-circle-powers}{Signed powers with occurrence-sensitive computation}
For a loop \(\ell:a=a\), define \(P_\ell:C\to(a=a)\) by
\[
\begin{aligned}
P_\ell(\mathsf z)&=\mathsf{refl}_a,&
P_\ell(\mathsf{pos}(0))&=\mathsf{refl}_a\cdot\ell,&
P_\ell(\mathsf{pos}(\mathsf Sk))&=P_\ell(\mathsf{pos}(k))\cdot\ell,\\
P_\ell(\mathsf{neg}(0))&=\mathsf{refl}_a\cdot\ell^{-1},&
P_\ell(\mathsf{neg}(\mathsf Sk))&=P_\ell(\mathsf{neg}(k))\cdot\ell^{-1}.
\end{aligned}
\]
Write \(\ell^n=P_\ell(G(n))\). In particular
\(\ell^0\equiv_{\mathrm j}\mathsf{refl}_a\), whereas
\(\ell^1=\ell\) is generally only propositional. These are the source's
positive, zero and negative powers, with its two signed recursions made
explicit \cite[\S6.10; \S8.1]{HoTT2013PublicReleaseFoundations}.
\end{hbdefinition}

\begin{hblemma}{thm:p07-circle-power-coherence}{proved=annex:p07-circle-power-coherence}{The predecessor coherence is a proved path}
For every \(n:\mathbb Z\) and loop \(\ell\),
\[
\ell^{\mathsf{succ}(n)}=\ell^n\cdot\ell,\qquad
\ell^{\mathsf{pred}(n)}=\ell^n\cdot\ell^{-1},\qquad
\ell^{\mathsf{pred}(n)}\cdot\ell=\ell^n.
\]
These are propositional identities, including at the two sign boundaries.
\end{hblemma}

\begin{hbdefinition}{def:p07-circle-code}{The univalent code family}
Circle recursion defines \(\mathrm{code}:S^1\to\mathcal U\) by
\[
\mathrm{code}(\mathsf{base})\equiv_{\mathrm j}\mathbb Z,\qquad
\operatorname{ap}_{\mathrm{code}}(\mathsf{loop})
 =\operatorname{ua}(e_{\mathsf{succ}}).
\]
The second equation is the propositional path computation of circle recursion.
The nondependent codomain \(\mathcal U\) does not turn it into a judgmental rule.
\end{hbdefinition}

\begin{hblemma}{thm:p07-circle-monodromy}{proved=annex:p07-circle-monodromy}{Transport is successor and inverse transport is predecessor}
For \(n:\mathbb Z\),
\[
T_{\mathrm{code}}(\mathsf{loop},n)=\mathsf{succ}(n),\qquad
T_{\mathrm{code}}(\mathsf{loop}^{-1},n)=\mathsf{pred}(n).
\]
\end{hblemma}

\begin{hbdefinition}{def:p07-circle-encode}{Encoding by transport}
For \(x:S^1\) and \(p:\mathsf{base}=x\), put
\[
\mathrm{encode}_x(p)=T_{\mathrm{code}}(p,0).
\]
In particular \(\mathrm{encode}_{\mathsf{base}}(\mathsf{refl})
\equiv_{\mathrm j}0\).
\end{hbdefinition}

\begin{hblemma}{thm:p07-circle-decoder}{proved=annex:p07-circle-decoder}{A decoder with its actual circle-induction lift}
There is a dependent function
\[
\mathrm{decode}:\prod_{x:S^1}
 \bigl(\mathrm{code}(x)\to(\mathsf{base}=x)\bigr)
\]
with \(\mathrm{decode}_{\mathsf{base}}(n)
\equiv_{\mathrm j}\mathsf{loop}^{n}\).
Its supplied loop lift is the function path whose component at \(n\) is
\[
T_{x\mapsto(\mathrm{code}(x)\to(\mathsf{base}=x))}
   (\mathsf{loop},\mathsf{loop}^{(-)})(n)
 =\mathsf{loop}^{\mathsf{pred}(n)}\cdot\mathsf{loop}
 =\mathsf{loop}^{n}.
\]
The first comparison uses function transport and inverse monodromy; the
second uses the proved signed coherence, not an assumed decoder.
\end{hblemma}

\begin{hblemma}{thm:p07-circle-winding}{proved=annex:p07-circle-winding}{Signed transport and the winding calculation}
For every \(m,n:\mathbb Z\),
\[
T_{\mathrm{code}}(\mathsf{loop}^{n},m)=m+n,\qquad
\mathrm{encode}_{\mathsf{base}}(\mathsf{loop}^{n})=n.
\]
The zero case is transport reflexivity; the positive cases iterate successor
and the negative cases iterate predecessor.
\end{hblemma}

\begin{hbtheorem}{thm:p07-circle-decode-encode}{proved=annex:p07-circle-decode-encode}{The first inverse homotopy}
For \(x:S^1\) and \(p:\mathsf{base}=x\),
\[
\eta_x(p):\mathrm{decode}_x(\mathrm{encode}_x(p))=p.
\]
This is based identity elimination with a variable endpoint \(x\), not
identity elimination on a fixed loop type.
\end{hbtheorem}

\begin{hbtheorem}{thm:p07-circle-encode-decode}{proved=annex:p07-circle-encode-decode}{The second inverse homotopy with circle coherence}
For \(x:S^1\) and \(c:\mathrm{code}(x)\),
\[
\epsilon_x(c):\mathrm{encode}_x(\mathrm{decode}_x(c))=c.
\]
The circle-induction motive is a family of function types of identity proofs,
and its loop lift is supplied explicitly in the annex.
\end{hbtheorem}

\begin{hbtheorem}{thm:p07-circle-loop-equivalence}{proved=annex:p07-circle-loop-equivalence}{The loop space is the integers}
The encoder and decoder give, for every \(x:S^1\),
\[
(\mathsf{base}=x)\simeq\mathrm{code}(x);
\quad\text{in particular}\quad
\Omega(S^1,\mathsf{base})\simeq\mathbb Z.
\]
The loop space is a set. The pointed comparison sends reflexivity to \(0\),
and \(\mathrm{encode}(p\cdot q)=\mathrm{encode}(p)+\mathrm{encode}(q)\).
Thus it identifies loop concatenation with the source's integer addition.
\end{hbtheorem}

\begin{hbcorollary}{thm:p07-circle-cover-contractible}{proved=annex:p07-circle-cover-contractible}{The total code space is contractible}
The type \(\sum_{x:S^1}\mathrm{code}(x)\) is contractible with center
\((\mathsf{base},0)\).
\end{hbcorollary}

\paragraph{Worked paths and limits of computation.}
Here \(1=E(\mathsf{pos}(0))\), \(2=E(\mathsf{pos}(1))\),
\(-1=E(\mathsf{neg}(0))\), and \(-2=E(\mathsf{neg}(1))\);
the indices \(1\) in the tags are natural-number successors of zero.
The loops \(\mathsf{loop}^{2}\) and \(\mathsf{loop}^{-2}\) encode as \(2\)
and \(-2\): their transports are respectively
\(\mathsf{succ}(\mathsf{succ}(0))\) and
\(\mathsf{pred}(\mathsf{pred}(0))\).
Transport along \(\mathsf{loop}^{-1}\cdot\mathsf{loop}\) first subtracts
one and then adds one, giving \(0\). Its equality with reflexivity is the
path inverse law; it is not a new syntactic reduction rule. Decoding \(-1\)
gives \(\mathsf{refl}\cdot\mathsf{loop}^{-1}\), propositionally
\(\mathsf{loop}^{-1}\). If a path
\(\mathsf{loop}=\mathsf{refl}\) existed, applying the encoder and monodromy
would identify \(1\) and \(0\), contradicting the disjoint canonical integer
tags proved in the annex. In contrast, the encoder does not assert that
all proofs of the univalence axiom are judgmentally equal.

\begin{proofboundary}{def:p07-circle-native-scope}{Native scope and distinct proof routes}
The two inverse homotopies above complete the main encode--decode route of
the June 2013 \S8.1. The source also gives a second route using the general
flattening lemma and a third route using the identity-system eliminator.
The derived total-space contraction here is a consequence of the two
inverse homotopies. The separate coequalizer flattening and identity-system
routes, with their own dependent path calculations, are supplied in
Section~\ref{sec:p07-hott-identity-systems-and-flattening}; they are not
inferred merely from those two inverse homotopies.
The source's statements about \(\pi_1\) and higher \(\pi_n\) additionally
use its truncation and iterated-loop interfaces; no general sphere
calculation, general homotopy-group functor, or completion of all of
Chapters 6 and 8 is inferred from this circle equivalence.
Without univalence the particular universe loop defining \(\mathrm{code}\)
has not been supplied; without the HIT path computation or funext the
decoder lift above has not been obtained. A merely typed decoder function,
a citation of the winding theorem, or fixed-endpoint path induction would
not discharge these obligations.
\end{proofboundary}
